\section{Navigation}

A space is only usable as a computer if you can find your way around it, and the mesh has a fast address for every
cell. Finding any one cell among astronomically many costs almost nothing.

Because the layer growth follows a fixed recurrence, the same recurrence that gave the warp factor, every dock has a
short canonical name, a string of digits read off its position in the growth tree, and a step to a neighbor is plain
arithmetic on that string. So a position in the mesh becomes a number, and moving to an adjacent cell becomes adding
to that number, with no search required.

The payoff is logarithmic cost. The number of cells within a given distance grows exponentially, so the distance to any
cell is only about the logarithm of the total number of cells, and a name is about that many digits long. Finding,
addressing, or routing to any cell therefore costs about the logarithm of the number of cells, which is fast: among a
billion cells, a route is a few dozen steps. An experiment measures this, finding that navigation and neighbor-stepping
run in logarithmic cost on the curved mesh, far below the cost on a flat lattice of the same size, where reaching
across the space costs a power of the size rather than its logarithm.

This is the practical face of the hyperbolic geometry. The same exponential branching that crowds information onto the
boundary and gives holography also makes the space shallow, every cell only a logarithmic number of steps from any other
through the branching interior. So the mesh is not a vast flat expanse you must trudge across, it is a deep tree you
can reach into, which is what lets a bounded region coordinate a huge interior, and what a mind needs to hold and search
a world.
